
\mode<presentation>
{
 \usetheme{default}      % or try Darmstadt, Madrid, Warsaw, ... 
 %\usecolortheme{default} % or try albatross, beaver, crane, ...
 \usecolortheme{TUNI} % mjh, 2018
%  \usefonttheme{default}  % or try serif,
%  \usefonttheme{serif}  % or try serif,structurebold, ...
 \setbeamertemplate{navigation symbols}{}
 \setbeamertemplate{caption}[numbered]
} 

\usepackage[english]{babel}
%\usepackage[utf8]{inputenc}
\usepackage[OT1]{fontenc} %maybe this works better with \href{with ä}
\usepackage{mwe}
\usepackage[absolute,overlay]{textpos}
\usepackage{pgf}  
\usepackage{tikz}
\usepackage{calc}

\usepackage{mathtools}
\usepackage{animate}
\usepackage{transparent}
\usepackage{tikz}
\usepackage{stackengine}

\usepackage{threeparttable}
\usepackage[labelformat=empty,font=scriptsize,skip=0pt,
justification=justified,singlelinecheck=false]{caption}

\mathchardef\mhyphen="2D % Define a "math hyphen"
\newcommand{\micron}{{\textmu}m} %$\mu$m
% \newcommand{\Alert}[1]{\alert{\textbf{#1}}}
\newcommand{\Alert}[2][blue]{\textcolor{#1}{\textbf{#2}}}

% recursive evilness
%\renewcommand{\textbf}[2][black]{\textcolor{#1}{\textbf{#2}}}
 
\usepackage{enumitem}

\usepackage{tikz} 

%% PRESENTATION ASPECT RATIO AFFECTS FOLLOWING POSITIONERS:
%% CITATIONS %% 
% 16:9 aspect ratio
% \newcommand\BRciteXpos{159mm} 	% 127 mm
% \newcommand\BRciteYpos{88mm}	% 94 mm
% %% PAGE NUMBERING %%
% % 4:3 aspect ratio
% % \newcommand\framenumXpos{98mm} 	% 118 mm
% % 16:9 aspect ratio
% \newcommand\framenumXpos{130mm} 	% 118 mm

% \draw[step=0.5cm,gray,very thin] (-1.9,-1.9) grid (5.9,5.9); % for positioning stuff in drafting mode

\graphicspath{{Figures/}{Figures/Logos/}{Figures/misc/}{Figures/Prezi/}{Figures/Intro/}{Figures/MRES/}{Figures/NDDA/}{Figures/LessIsMore/}{Figures/3DNLO/}{Figures/Seminar2020/}{Figures/DeepLearning/}{Figures/highQ/}{Figures/Transients/}{Figures/PREIN_lukio_kuvat/}}
% Organize figures into folders: 
% Figures/NLO		(NLO schematics)
% Figures/misc 		(TUT logos etc.)
% Figures/Logos 	(TUT logos etc.)
% Figures/Prezi 	(this presentation figs)
% Figures/?			(other presentation figs)
% etc.

%% Automated titlecase:
\newcommand{\tc}[1]{\titlecap{#1}}

%%%%%%%%%%%%%% FOOTLINE SIZE: %%%%%%%%%%%%%%%%%
% \makeatletter
% \defbeamertemplate*{footline}{Dan P theme}
% {
%   \leavevmode%
%   \hbox{%
%   \begin{beamercolorbox}[wd=.2\paperwidth,ht=2.25ex,dp=1ex,center]{author in head/foot}%
%     \usebeamerfont{author in head/foot}\insertshortauthor\expandafter\beamer@ifempty\expandafter{\beamer@shortinstitute}{}{~~(\insertshortinstitute)}
%   \end{beamercolorbox}%
%   \begin{beamercolorbox}[wd=.6\paperwidth,ht=2.25ex,dp=1ex,center]{title in head/foot}%
%     \usebeamerfont{title in head/foot}\insertshorttitle
%   \end{beamercolorbox}%
%   \begin{beamercolorbox}[wd=.2\paperwidth,ht=2.25ex,dp=1ex,right]{date in head/foot}%
%     \usebeamerfont{date in head/foot}\insertshortdate{}\hspace*{2em}
% \insertframenumber{} / \inserttotalframenumber\hspace*{2ex} 
%   \end{beamercolorbox}}%
%   \vskip0pt%
% }
% \makeatother

%%%%%%%%%%%%%% NEW COMMANDS: %%%%%%%%%%%%%%%%%

% Itemize list -customization:
\setbeamercolor{itemize item}{fg=blue}  % mainBlue darkblue
\setbeamertemplate{itemize item}[circle]  % triangle
\newcommand\mjhLineSpace{-1mm}
\newcommand\FontThanks{\fontsize{22}{7.2}\selectfont}

% To highlight results using colors:
\usepackage{xcolor,soul} %allows you to use \hl{} to highlight text
\definecolor{magenta}{RGB}{230,0,230} %\definecolor{orange}{rgb}{1,0.5,0}
\newcommand{\magenta}[1]{{\color{magenta}#1}}  
\definecolor{green}{RGB}{0,191,0} %\definecolor{orange}{rgb}{1,0.5,0}
\newcommand{\green}[1]{{\color{green}#1}}  %% You can create your own 
\definecolor{expblue}{RGB}{0,0,255} %\definecolor{orange}{rgb}{1,0.5,0}
\newcommand{\blue}[1]{{\color{expblue}#1}}  %% You can create your own 

% %%% WIKI-forwarding \EMPH -command
% \DeclareTextFontCommand{\emphLink}{\em}%\bfseries\em
% \DeclareDocumentCommand\EMPH{o g}{%
%     {\IfNoValueT{#1}%
%     {\href{https://en.wikipedia.org/wiki/#2}{\emphLink{#2}}}%
%     \IfNoValueF {#1} {\href{#1}{\emphLink{#2}}}%
%     }%
% }

%%%%%%%%%%%%%%%%%%%% EMPH %%%%%%%%%%%%%%%%%%%%
\let\emph\relax % there's no \RedeclareTextFontCommand
% \DeclareTextFontCommand{\emph}{\bfseries\em}
%% make lmgtfy links to emph'd terms..
\DeclareTextFontCommand{\emph}{\bfseries\em}
\DeclareTextFontCommand{\emphLink}{\bfseries\em}

% \newcommand{\EMPH}[1]{\href{https://lmgtfy.com/?q=#1}{\emph{#1}}}
%https://lmgtfy.com/?q=index+of+refraction
%% wiki version:
% \newcommand{\EMPH}[1]{\href{https://en.wikipedia.org/wiki/#1}{\emph{#1}}}
% \newcommand{\EMPH}[1]{\href{https://en.wikipedia.org/wiki/#1}{\emph{#1}}}

% https://stackoverflow.com/questions/1812214/latex-optional-arguments
\usepackage{xparse}
% \DeclareDocumentCommand\EMPH{ m g }{%
%     {\IfNoValueT{#2}%
%     {\href{https://en.wikipedia.org/wiki/#1}{\emph{#1}}}%
%     \IfNoValueF {#2} {\href{#2}{\emph{#1}}}%
%     }%
% }

\usepackage{mfirstuc}
\makeatletter
\newcommand*\define[1]{%
  %\textit{#1}%
  \index{#1@\protect\capitalisewords{#1}}%
   }
\makeatother
 
\DeclareDocumentCommand\EMPH{ m g g}{%
    {\IfNoValueT{#3}%
    {\define{#1}}%
    \IfNoValueF {#3} {\define{#3}}%
    }%
    {\IfNoValueT{#2}%
    {\href{https://en.wikipedia.org/wiki/#1}{\emphLink{#1}}}%
    \IfNoValueF {#2} {\href{#2}{\emphLink{#1}}}%
    }%
}
% \newcommand{\example}[2][YYY]{Mandatory arg: #2;
%                                  Optional arg: #1.}
%% EMPH withouth indexing:                                
\DeclareDocumentCommand\EMPHn{ m g g}{%
    {\IfNoValueT{#3}%
    {}%
    \IfNoValueF {#3} {\define{#3}}%
    }%
    {\IfNoValueT{#2}%
    {\href{https://en.wikipedia.org/wiki/#1}{\emphLink{#1}}}%
    \IfNoValueF {#2} {\href{#2}{\emphLink{#1}}}%
    }%
}

%%%%%%%%%%%%%%%%%%%%%%%%%%%%%%
%% Make hyperlink to bottom left:
\newcommand{\hlinkBL}[2]
{\setlength{\TPHorizModule}{\textwidth}
 \setlength{\TPVertModule}{\textwidth}
 \begin{textblock}{0.5}(0.5,0.860)  % Left Bottom = (0.01,0.86)
 \hyperlink{#1}{\beamergotobutton{#2}}   
\end{textblock}}
%%%%%%%%%%%%%%%%%%%%%%%%%%%%%%
%% Make horizontal navigation bar:
\newcommand{\insertCentHorNavBar}
{\setbeamertemplate{headline}{%
\leavevmode%
  \hbox{%
    \begin{beamercolorbox}[wd=\paperwidth,ht=2.5ex,dp=1.125ex]{palette quaternary}%
  % \insertsectionnavigationhorizontal{\paperwidth}{}{\hskip0pt plus1filll}
  % \insertsectionnavigationhorizontal{\paperwidth}{\hskip0pt plus1filll}{\hskip0pt plus1filll}
    \insertsectionnavigationhorizontal{\paperwidth}{\hfill\hfill}{\hfill\hfill}
    \end{beamercolorbox}%
  }
}}

%%%%%%%%%%%%%%%%%%%%%%%%%%%%%%
%% INSERT TEXT TO CUSTOM POSITION:
\newcommand{\textToArbPos}[2]
{\setlength{\TPHorizModule}{\textwidth}
 \setlength{\TPVertModule}{\textwidth}
 \begin{textblock}{0.5}(#1){#2}  % Left Bottom = (0.01,0.86)
 \end{textblock}}

%%%%%%%%%%%%%%%%%%%%%%%%%%%%%%
%%% Customize colors, Fig. captions and footlines %%% 

%\setbeamertemplate{caption}{\raggedright\insertcaption\par} % removes Fig. 1 from the captions.
\setbeamertemplate{caption}{\centering\insertcaption\par} % removes Fig. 1 from the captions.
\setlength\abovecaptionskip{2pt}
% \setbeamertemplate{caption}{\centering\insertcaption\par} % removes Fig. 1 from the captions.
\setbeamercolor{footline}{bg=white,fg=darkblue}
\setbeamerfont{caption}{size=\scriptsize}

\usepackage{xcolor}
\newcommand{\highlight}[1]{%
  \colorbox{mainViolet}{$\displaystyle#1$}} % mainGreen!50

%%%%%%%%% SUB-SECTION AUTONAMING %%%%%%%%%%%%% 
%% extract frametitle to auto-name sub-sections
\addtobeamertemplate{frametitle}{
\subsection{\insertframetitle}}



%%%%%%%%%%%%%%%%%%%%%%%%%%%
%% define the golden ratio lengths:
\newcommand\gra{0.618}
\newcommand\TUTsbWdth{0.15cm}
%\newcommand{\TUTlogoW}[1]{%
%\number\numexpr#1+\TUTsbWdth\relax%
%}
%%%%% INSERT TUNI LOGO %%%%%
%https://tex.stackexchange.com/questions/53781/how-can-i-include-the-logo-in-some-slides-and-remove-in-others-using-beamer
% \newif\ifplacelogo % create a new conditional
% \placelogotrue % set it to true
% \ifplacelogo%
%%%%% INSERT TUNI LOGO %%%%%
%https://tex.stackexchange.com/questions/53781/how-can-i-include-the-logo-in-some-slides-and-remove-in-others-using-beamer
% \newif\ifplacelogo % create a new conditional
% \placelogotrue % set it to true
% \ifplacelogo%
% \setbeamertemplate{footline}[text line]
% {\parbox[c][1.4cm]{\linewidth}
% {\hspace*{-0.9cm}  
% %  \includegraphics[height=1cm]{Figures/Logos/TUNI_logo_violet.eps}}  % was 1.1cm earlier
%  \includegraphics[height=1cm]{logo_TAU_2rivi_fi_violetti_RGB}}  % was 1.1cm earlier
% }%

\newcommand{\nologo}{\setbeamertemplate{footline}[text line]{}}

%%%%% INSERT TUT LEFT-BAR %%%%%
\setlength{\TPHorizModule}{10mm}
\setlength{\TPVertModule}{\TPHorizModule}
\textblockorigin{0mm}{1\paperheight} % start everything near the top-left corner
% \setlength{\parindent}{0pt}

%%%%% CUSTOMIZE TABLE OF CONTENT (ToC) STYLE %%%%%
% SECTION
\setbeamertemplate{section in toc}[circle]
\setbeamercolor{section in toc}{fg=darkblue}
\setbeamertemplate{section in toc}
{{\color{darkblue}\inserttocsectionnumber.}~\inserttocsection}
% SUB-SECTION
\setbeamercolor{subsection in toc}{fg=black}
\setbeamerfont{subsection in toc}{size=\footnotesize} % \small 

%\setbeamertemplate{section in toc}{\inserttocsectionnumber.~\inserttocsection}
%%%%%%%%%%%%%%%%%%%%%%%%%%%%%%%%%%%%%%%%%%%%%%%%%%

 
